\scanpage{29}
\[
\begin{aligned}
C(\omega_0^u)^{u^{-1}}={}&\left(\unsure{\varpi_0|\bar\omega_0}+\frac{\omega_0}{1+\sum\alpha_i\omega_0^i}\unsure{\varpi_1|\bar\omega_1}+\left(\frac{\omega_0}{1+\sum\alpha_i\omega_0^i}\right)^2\unsure{\varpi_2|\bar\omega_2}\right)\\
&\cdot\left[1+\sum(i+1)\alpha_i\left(\frac{\omega_0}{1+\sum\alpha_j\omega_0^j}\right)^{\unsure{i}}\right.\unsure{\big]}\\
&+\sum\left(\frac{\omega_0}{\unsure{\sum1+\alpha_j\omega_0^j|1+\sum\alpha_j\omega_0^j}}\right)^{i+1}\otimes d\alpha_i.
\end{aligned}
\]
Supposons $\alpha_2=0$, et calculons mod $\Omega^4$
\[
\begin{aligned}
C(\omega_0^u)^{u^{-1}}
&\underset{\substack{\mathrm{mod}\\\sum_{i\geq4}\Omega^i\otimes\Omega^1_{X/S}}}{\equiv}
(\unsure{\varpi_0|\bar\omega_0}+\omega_0\unsure{\varpi_1|\bar\omega_1}+\omega_0^2\unsure{\varpi_2|\bar\omega_2})
[1+\deleted{\unsure{(i+1)\alpha_i}}\inserted{4\alpha_3\omega_0^3}]\\
&\unsure{+}\\
&=\deleted{\unsure{\varpi_2}}C(\omega_0)+4\alpha_3\omega_0^3\unsure{\varpi_0|\bar\omega_0}
\end{aligned}
\]
\margin{il faut $\alpha_3=0$}
Supposons $\alpha_2=0\quad\unsure{\alpha_{n-1}|\alpha_n}=0$ $(n\geq3)$ calculons mod $\Omega^{n+1}$
\inserted{$\alpha_n=?$}
\[
\begin{aligned}
C(\omega_0^u)^{u^{-1}}&=(\unsure{\varpi_0|\bar\omega_0}+\omega_0\unsure{\varpi_1|\bar\omega_1}+\omega_0^2\unsure{\varpi_2|\bar\omega_2})(1+(n+1)\alpha_n\omega_0^n)\\
&=\deleted{\unsure{\varpi_0}}C(\deleted{\unsure{\omega}}\omega_0)+\deleted{\unsure{\varpi_0\lambda}}\underline{(n+1)\alpha_n\omega_0^n}\unsure{\varpi_0|\bar\omega_0}
\end{aligned}
\]
donc il faut $(n+1)\alpha_n=0$.

Calculons mod $\Omega^3$ en \inserted{ne} \unsure{supposant plus} $\alpha_2$,
Je trouve
\[
\begin{aligned}
C(\omega_0^u)^{u^{-1}}&=(\unsure{\varpi_0|\bar\omega_0}+\omega_0\unsure{\varpi_1|\bar\omega_1}+\omega_0^2\unsure{\varpi_2|\bar\omega_2})(1+3\alpha_2\omega_0^2)\\
&=C(\omega_0)+\underline{3\alpha_2}\omega_0^2\underline{\unsure{\varpi_0|\bar\omega_0}}
\end{aligned}
\]
\uline{Prop.} Soit \inserted{$\operatorname{Conn}(P)$} l'ens des connexions \inserted{$c$} de $P$ (ou de l'Alg. \inserted{de Lie} $(\unsure{H|\mathfrak g})$), \deleted{\unsure{soit}} $G_0$ \deleted{l'ens} groupe des autom $u$ de $\widehat S$ \inserted{(alg. augm.)} (qui sont l'identité mod $\unsure{\mathcal J^4|\mathcal J^3}$), considérons l'application
\diagram{\begin{tikzcd}
(c,u) \arrow[r,maps to] & \deleted{\unsure{\widehat c}}\deleted{\unsure{u\widehat c}}u(\widehat c)\\
\operatorname{Conn}P\inserted{\times G_0} \arrow[r] & \operatorname{Conn}_{\unsure{\Omega_0|\omega_0}}(\widehat S)\deleted{\unsure{\mathrm{mod}}}
\end{tikzcd}}
\figure{\begin{tikzpicture}\draw (0,0)--(0,-2.7);\end{tikzpicture}}
\uline{Cette application est bijective} ($\unsure{X|S}$ de car. $0$)
% archive foliation: 28
